% !PRINTED 469
\weberpage{469}

\[
\sigma_a = (\varrho, \varrho_a), \qquad a = 0, 1, 2 \ldots \mu - 1,
\]
where by $\sigma_0$ or $\sigma$ is understood the so-called \gesperrt{identical substitution} $(\varrho, \varrho)$, which consists in $\varrho$ being replaced by itself, so that all numbers in $\Omega(\varrho)$ remain unchanged.

If in any one of the roots (3)
\[
\varrho_h = \Theta_h(\varrho)
\]
we carry out the substitution $\sigma_a$, then $\varrho_h$ passes into another root $\varrho_k$, which is determined by
\[
\varrho_k = \Theta_h(\varrho_a) = \Theta_h \Theta_a(\varrho) = \Theta_k(\varrho).
\]
Thus $\varrho_k$ is the same function of $\varrho_a$ as $\varrho_h$ is of $\varrho$.

An arbitrary quantity $\omega = \varphi(\varrho)$ of the field $\Omega(\varrho)$ passes, by the substitution $\sigma_a$, into $\omega_a = \varphi(\varrho_a)$. If we express $\omega$ in terms of $\varrho_h$ and $\omega_a$ in terms of $\varrho_k$:
\[
\omega = \psi(\varrho_h), \qquad \omega_a = \psi(\varrho_k),
\]
then one sees that $\omega$ passes into $\omega_a$ by the substitution
\weq{9}{\sigma_a = (\varrho_h, \varrho_k) = (\varrho_h, \Theta_h \Theta_a(\varrho)),}
and this is only another expression for the substitution $(\varrho, \varrho_a)$.

Here $\varrho_h$ is an arbitrary root of $g(t)$, and the associated $\varrho_k$ is determined by $\sigma_a$. But for a given $\sigma_a$ one may also take $\varrho_k$ arbitrarily and determine the associated $\varrho_h$, by letting the index $h$ in $\Theta_h(\varrho_a)$ run through the values $0, 1, 2 \ldots \mu - 1$, whereby each root $\varrho_k$ makes its appearance once. Hence:

\gesperrt{3. Every substitution $\sigma_a$ can be represented in the form $(\varrho_h, \varrho_k)$, in which either $\varrho_h$ or $\varrho_k$ is an arbitrarily given one of the $\mu$ roots $\varrho$, while the other is completely determined by $\sigma_a$.}

% [sic] The German source breaks off grammatically after „in die gleiche“; the implied completion is „Grösse übergehen“ (pass into the same quantity). The English supplies the implied noun while the German transcription remains diplomatic.
The number of all mutually distinct substitutions $\sigma_a$ is therefore, the identical substitution included, equal to the degree of the field $\Omega(\varrho)$, which we have already denoted above by $\mu$. Each of these substitutions carries the totality of the quantities of the field $\Omega(\varrho)$ into itself, in such a way that each passes into a definite other one, and never two distinct quantities into the same one [formulae (7), (8)].

We therefore call the $\sigma_a$ the \gesperrt{substitutions of the field} $\Omega(\varrho)$.

If $\omega = \varphi(\varrho)$ remains unchanged in the field $\Omega(\varrho)$ when $\varrho$ is replaced by $\varrho_a$, that is, if
%CONT>>
